\section{A toy imperative language}
\label{sec:syntax}

Programs act on a memory holding one natural number per variable. There is no
allocation, no failure and no procedure call in the core language; procedures
are discussed in the companion note.

\begin{definition}[stores]
\label{def:state}
Fix a countable set $\Var$ of \emph{variables}. A \emph{store} is a total map
$\sigma : \Var \to \Nat$. We write $\Stores$ for the set of stores. Totality is
a convenience, not a commitment: it removes the ``uninitialised variable'' case
from every statement below.
\end{definition}

\begin{definition}[arithmetic expressions]
\label{def:aexp}
Arithmetic expressions are generated by
\[
  a \;::=\; n \;\mid\; x \;\mid\; a_1 + a_2,
  \qquad n \in \Nat,\; x \in \Var .
\]
The \emph{free variables} $\fv(a) \subseteq \Var$ of an expression are those
occurring in it, read off the syntax.
\end{definition}

\begin{definition}[boolean expressions]
\label{def:bexp}
Boolean expressions are generated by
\[
  b \;::=\; \kw{true} \;\mid\; a_1 \le a_2 \;\mid\; \lnot b \;\mid\; b_1 \land b_2 ,
\]
where $a_1, a_2$ range over the arithmetic expressions of \Cref{def:aexp}.
\end{definition}

\begin{definition}[commands]
\label{def:cmd}
Commands are generated by
\[
  c \;::=\; \skipk
     \;\mid\; x := a
     \;\mid\; c_1 ; c_2
     \;\mid\; \ifk\ b\ \thenk\ c_1\ \elsek\ c_2
     \;\mid\; \whilek\ b\ \dok\ c ,
\]
with $a$ as in \Cref{def:aexp} and $b$ as in \Cref{def:bexp}. Nothing in the
language can diverge except $\whilek$, and nothing can fail at all; this is
what makes partial correctness the only interesting judgement.
\end{definition}
